\chapter{Conclusioni e sviluppi futuri}

Questa tesi ha presentato la progettazione e l’implementazione di un tool per la scrittura e l’esecuzione di programmi nel linguaggio \textit{Racing Choreographies}.

Come discusso nel Capitolo 2, la macchina astratta implementa una versione semplificata della semantica operazionale rispetto a quella descritta nel modello teorico del linguaggio \textit{Racing Choreographies}.

Il contributo principale di questa tesi consiste nell’aver realizzato uno strumento che costituisce una prima base operativa di un linguaggio ancora in fase sperimentale nonostante non copra ancora tutti gli aspetti teorici e semantici presenti nella specifica completa.

Proprio per questo il lavoro svolto apre a diverse direzioni di sviluppo futuro. Un primo sviluppo naturale consiste nell’estendere il simulatore in modo da che realizzi la semantica completa definita nel modello teorico del linguaggio.

Inoltre il modello teorico del linguaggio considera ulteriori proprietà che non sono state implementate in questo lavoro tra cui: le condizioni di well-formedness delle coreografie, il supporto al massimo parallelismo tramite regole di delay basate sulle condizioni di Bernstein e le proprietà strutturali necessarie alla corretta endpoint projection.

Un’ulteriore direzione di lavoro, particolarmente importante nel contesto della programmazione coreografica, è l’introduzione della \textit{Endpoint Projection}. Il tool sviluppato in questa tesi opera interamente a livello di coreografia globale ma un passo successivo sarà proiettare tale descrizione sui singoli endpoint ottenendo comportamenti locali coerenti con la specifica globale.

In conclusione, il risultato principale di questa tesi è aver trasformato una specifica formale sperimentale in un primo artefatto eseguibile. Il tool realizzato offre una base utile per futuri studi sul linguaggio \textit{Racing Choreographies} e per successive evoluzioni della sua implementazione.